% answers_short.tex —— 2006 年·普通高等学校招生考试·重庆卷文科 · 简短答案速查
% 由 tools/make_short_answers.py 自 answers.tex 的 \gksolhead/\gkanswer 提取，
% 并与 2006/2006_数学真题_重庆文_答案要点.md 逐题比对后生成，勿手改（改源 MD 后重跑）。

\begin{gknotice}
\textbf{使用说明}：本节只给最终答案便于快速核对，完整推导见解析版。答案为编写者独立推导，\textbf{非官方标准答案}。
\end{gknotice}

\gkpart{选择题}
\begin{center}\begin{tabular}{ccccccccccccc}
\toprule
题号 & 1 & 2 & 3 & 4 & 5 & 6 & 7 & 8 & 9 & 10 & 11 & 12 \\
\midrule
答案 & D & D & C & D & B & C & C & D & B & B & A & A \\
\bottomrule
\end{tabular}\end{center}

\gkpart{填空题}
\begin{itemize}[leftmargin=2.6em]
  \item[\textbf{13.}] $-$2。
  \item[\textbf{14.}] 2n$-$1。
  \item[\textbf{15.}] (2,+∞)。
  \item[\textbf{16.}] a>1/2。
\end{itemize}

\gkpart{解答题}
\begin{itemize}[leftmargin=2.6em]
  \item[\textbf{17.}] (1) 1/6；(2) 5/72
  \item[\textbf{18.}] (1) $\omega=\frac12$；(2) a=(√3+1)/2
  \item[\textbf{19.}] (1) a=1，b=$-$3；(2) 递增区间为 ($-$∞,$-$1] 与 [3,+∞)，递减区间为 [$-$1,3]
  \item[\textbf{20.}] (1) $\frac{\pi}{6}$（即 $\arccos\left(\frac{\sqrt3}{2}\right)$）；(2) 2
  \item[\textbf{21.}] (1) a=2，b=1；(2) k<$-$1/3
  \item[\textbf{22.}] (1) $x_ns_n=-4$；(2) 题面等式成立
\end{itemize}
